%%%%%%%%%%%%%%%%%%%%%%%%%%%%%%%%%
%
%       EXERCÍCIO
%
%%%%%%%%%%%%%%%%%%%%%%%%%%%%%%%%%

\ifdefstring{\atividade}{prova}{%
    \ifdefstring{\modo}{objetivo}{%
        \renewcommand{\valorquestao}{\ValorQObj\ ponto}
    }{%
        \renewcommand{\valorquestao}{\ValorQDisc\ pontos}
    }%
}%

\begin{exercicioBanco}[\valorquestao]
Considere a função afim
%
\[
    \begin{array}{rcl}
        f: \, \bbR & \to & \bbR \\
        x & \mapsto & \dfrac{3}{4}x - \dfrac{1}{2}.
    \end{array}
\]
%
Calcule \(f(x)\), para \(x \in \left\{-\dfrac{4}{3}, -\dfrac{2}{3}, 0, \dfrac{5}{2}\right\}\).
%

% Define as alternativas
\newcommand{\alternativas}{%
    \begin{center}
        \begin{tabularx}{\textwidth}{XXX}
            (a) \(\left\{-\dfrac{3}{2}, -1, \dfrac{1}{2}, \dfrac{11}{8}\right\}\). &
            (b) \(\left\{-\dfrac{3}{2}, -1, -\dfrac{1}{2}, \dfrac{11}{8}\right\}\). &
            (c) \(\left\{-\dfrac{3}{2}, -\dfrac{1}{2}, -\dfrac{1}{2}, \dfrac{11}{8}\right\}\). \\[15pt]
            (d) \(\left\{-1, -1, -\dfrac{1}{2}, \dfrac{11}{8}\right\}\). &
            (e) \(\left\{-\dfrac{3}{2}, -1, -\dfrac{1}{2}, \dfrac{5}{2}\right\}\).
        \end{tabularx}
    \end{center}
}

% Resposta correta
\newcommand{\resposta}{B}

% Lógica condicional para exibição
\ifdefstring{\atividade}{lista}{%
    \alternativas
    \vspace{0.5em}
    
    \noindent\textbf{Resposta:} letra \textbf{\resposta}.
}{%
    \ifdefstring{\modo}{objetivo}{%
        \alternativas
    }{%
        % modo = discursiva → não mostra alternativas
    }
}
\end{exercicioBanco}

